Starlink és una constel·lació de satèl·lits que ofereix serveis d'Internet de banda ampla utilitzant òrbites terrestres baixes per a fer més ràpida la comunicació. Cada satèl·lit té una massa de 227~kg i descriu una òrbita circular a una altura de 550~km per sobre de la superfície terrestre.

\begin{apartat}{1,25}
A partir de la llei de gravitació universal, deduïu l'expressió de la velocitat orbital en funció del radi orbital i de la massa de la Terra. A continuació, calculeu la velocitat orbital dels satèl·lits a l'altura de 550~km i determineu quantes voltes completes a la Terra descriu cada satèl·lit en un dia.
\end{apartat}

\begin{apartat}{1,25}
Per a evitar possibles col·lisions amb altres objectes, cada satèl·lit disposa d'un sistema de maniobra que, en detectar un risc, redueix l'altura de l'òrbita. En una d'aquestes maniobres, un satèl·lit baixa gradualment fins a una altura de 530~km. Deduïu l'expressió de l'energia mecànica en funció del radi orbital, la massa de la Terra i la massa del satèl·lit. A continuació, calculeu l'energia mecànica del satèl·lit en aquesta nova òrbita i justifiqueu per què l'energia mecànica orbital ha de ser negativa.
\end{apartat}

\dades{%
  $G = 6,67\times 10^{-11}\ \mathrm{N\,m^2\,kg^{-2}}$.\\
  Massa de la Terra: $M_\mathrm{T} = 5,98\times 10^{24}\ \mathrm{kg}$.\\
  Radi de la Terra: $R_\mathrm{T} = 6\,370\ \mathrm{km}$.}
